\QHeading{第1题　在\(\square\)里填上合适的数。}
\par\noindent


\PlainItem{\(63+88+37=88+(\){\color{answerred}63}\(+\){\color{answerred}37}\()\)}
\par

\PlainItem{\(159+72+41=72+(\){\color{answerred}159}\(+\){\color{answerred}41}\()\)}
\PlainItem{\(245+138+55=138+(\){\color{answerred}245}\(+\){\color{answerred}55}\()\)}
\PlainItem{\(76+208+24=208+(\){\color{answerred}76}\(+\){\color{answerred}24}\()\)}
\PlainItem{\(327+96+73=96+(\){\color{answerred}327}\(+\){\color{answerred}73}\()\)}
\PlainItem{\(184+57+16=57+(\){\color{answerred}184}\(+\){\color{answerred}16}\()\)}
\par\noindent

\QHeading{第2题　用简便方法计算下面各题。}

\CalcBlock[方法一：简便计算]{1}{\begin{aligned}[t]&\phantom{={}}\color{black}136+278+64\\&\color{answerred}=(136+64)+278\\&\color{answerred}=200+278\\&\color{answerred}=478\end{aligned}}
\par
\MethodBlock{方法二：按运算顺序直接计算（对照）}{本法只作结果核对；本题要求简便计算，前面的简便解法仍保留。}{\begin{aligned}[t]&\phantom{={}}136+278+64\\
&=414+64\\
&=478\end{aligned}}


\CalcBlock[方法一：简便计算]{2}{\begin{aligned}[t]&\phantom{={}}\color{black}(159+45)+55\\&\color{answerred}=159+(45+55)\\&\color{answerred}=159+100\\&\color{answerred}=259\end{aligned}}
\MethodBlock{方法二：按运算顺序直接计算（对照）}{本法只作结果核对；本题要求简便计算，前面的简便解法仍保留。}{\begin{aligned}[t]&\phantom{={}}(159+45)+55\\
&=204+55\\
&=259\end{aligned}}

\CalcBlock[方法一：简便计算]{3}{\begin{aligned}[t]&\phantom{={}}\color{black}248+175+52+25\\&\color{answerred}=(248+52)+(175+25)\\&\color{answerred}=300+200\\&\color{answerred}=500\end{aligned}}
\MethodBlock{方法二：交换、结合后先凑整}{保留每个数原来的加减号，先把能凑整的部分放在一起。}{\begin{aligned}[t]&\phantom{={}}248+175+52+25\\
&=(248+52)+175+25\\
&=300+175+25\\
&=475+25\\
&=500\end{aligned}}

\MethodBlock{方法三：按运算顺序直接计算（对照）}{本法只作结果核对；本题要求简便计算，前面的简便解法仍保留。}{\begin{aligned}[t]&\phantom{={}}248+175+52+25\\
&=423+52+25\\
&=475+25\\
&=500\end{aligned}}

\CalcBlock[方法一：简便计算]{4}{\begin{aligned}[t]&\phantom{={}}\color{black}317+86+83+14\\&\color{answerred}=(317+83)+(86+14)\\&\color{answerred}=400+100\\&\color{answerred}=500\end{aligned}}
\MethodBlock{方法二：交换、结合后先凑整}{保留每个数原来的加减号，先把能凑整的部分放在一起。}{\begin{aligned}[t]&\phantom{={}}317+86+83+14\\
&=(86+14)+317+83\\
&=100+317+83\\
&=417+83\\
&=500\end{aligned}}

\MethodBlock{方法三：按运算顺序直接计算（对照）}{本法只作结果核对；本题要求简便计算，前面的简便解法仍保留。}{\begin{aligned}[t]&\phantom{={}}317+86+83+14\\
&=403+83+14\\
&=486+14\\
&=500\end{aligned}}

\CalcBlock[方法一：简便计算]{5}{\begin{aligned}[t]&\phantom{={}}\color{black}125+78+75\\&\color{answerred}=(125+75)+78\\&\color{answerred}=200+78\\&\color{answerred}=278\end{aligned}}
\MethodBlock{方法二：按运算顺序直接计算（对照）}{本法只作结果核对；本题要求简便计算，前面的简便解法仍保留。}{\begin{aligned}[t]&\phantom{={}}125+78+75\\
&=203+75\\
&=278\end{aligned}}

\CalcBlock[方法一：简便计算]{6}{\begin{aligned}[t]&\phantom{={}}\color{black}208+182+1792\\&\color{answerred}=(208+1792)+182\\&\color{answerred}=2000+182\\&\color{answerred}=2182\end{aligned}}
\MethodBlock{方法二：按运算顺序直接计算（对照）}{本法只作结果核对；本题要求简便计算，前面的简便解法仍保留。}{\begin{aligned}[t]&\phantom{={}}208+182+1792\\
&=390+1792\\
&=2182\end{aligned}}

\par\noindent

\QHeading{第3题　解决问题。}

\Problem{①食堂运来大米900千克，第一周吃掉265千克，第二周吃掉235千克。食堂还剩多少千克大米？}
\par


\MethodTitle{方法一}
\Answer{\Formula{\begin{aligned}[t]&\phantom{={}}900-(265+235)\\&=900-500\\&=400\text{（千克）}\end{aligned}}\par 答：食堂还剩400千克大米。}

\MethodTitle{方法二：每用去一部分就求一次剩余}
\Answer{先减第一周吃掉的，再减第二周吃掉的。}
\Answer{\Formula{\begin{aligned}[t]&\phantom{={}}900-265-235\\
&=635-235\\
&=400\text{（千克）}\end{aligned}}}
\Answer{答：食堂还剩400千克大米。}
\Problem{②周末妈妈去超市购物，买零食花了132元，买果蔬花了208元，买生活用品花了68元。妈妈一共花了多少钱？}

\MethodTitle{方法一}
\Answer{\Formula{\begin{aligned}[t]&\phantom{={}}132+208+68\\&=(132+68)+208\\&=200+208\\&=408\text{（元）}\end{aligned}}\par 答：妈妈一共花了408元。}

\MethodTitle{方法二：按先后顺序逐次相加}
\Answer{先合计前两部分，再加第三部分；这是与凑整法对照的常规算法。}
\Answer{\Formula{\begin{aligned}[t]&\phantom{={}}132+208+68\\
&=340+68\\
&=408\text{（元）}\end{aligned}}}
\Answer{答：一共花了408元。}
